\documentclass[10pt,a4paper]{article}
\usepackage{cv-style}
\hypersetup{
  pdftitle={유리 허샤 - 상세 경력기술서},
  pdfauthor={유리 허샤},
  pdfsubject={로보틱스 소프트웨어 및 피지컬 AI},
  pdflang={ko-KR}
}

\begin{document}

\newgeometry{margin=15mm}
\cvheader
  {유리 허샤}
  {시니어 로보틱스 소프트웨어 엔지니어}
  {대한민국 서울}
  {\href{mailto:yurirocha15@gmail.com}{yurirocha15@gmail.com}\enspace\textbullet\enspace
   \href{https://www.yurirocha.com}{yurirocha.com}\enspace\textbullet\enspace
   \href{https://www.linkedin.com/in/yurirocha15/}{LinkedIn}\enspace\textbullet\enspace
   \href{https://github.com/yurirocha15}{GitHub}}

\cvmainsection{관심 분야}
{\small 로보틱스 \textbullet\ LLM 및 파운데이션 모델 \textbullet\ 강화학습 \textbullet\ 딥러닝 \textbullet\ DevOps \textbullet\ 소프트웨어 개발\par}

\cvmainsection{기술}
\cvproject{소프트웨어}{C++, Python, Go, Linux, JavaScript, Node.js, 소프트웨어 아키텍처}
\cvproject{실시간 시스템}{Real-time Linux, PREEMPT\_RT, Xenomai}
\cvproject{로보틱스}{ROS / ROS 2, MoveIt, OMPL, 모션 플래닝}
\cvproject{머신러닝}{PyTorch, ONNX}
\cvproject{시뮬레이션}{Isaac Sim, MuJoCo, Unity3D, Gazebo}
\cvproject{인프라}{Kubernetes, Docker, GitHub Actions}

\cvmainsection{경력}
\cvrole{두산로보틱스}{2025/08--현재}{시니어 소프트웨어 엔지니어, AI \& Software}{}
\begin{cvbullets}
  \item 고성능 로봇 컨트롤러의 핵심 실시간 소프트웨어를 설계·개발하고, 핵심 시스템의 태스크 관리 및 데이터 흐름 아키텍처를 설계하고 있습니다.
  \item 에이전틱 워크플로가 실제 로봇 하드웨어와 상호작용하고 이를 구동할 수 있도록 소프트웨어 인터페이스와 하네스를 개발하고 있습니다.
  \item 지속적인 개발과 배포를 지원하는 내부 Kubernetes 클러스터를 운영하고 CI/CD 및 LLMOps 파이프라인을 유지보수하고 있습니다.
  \item 개발자가 Kubernetes에 LLM과 시뮬레이션 환경을 손쉽게 배포하고 관련 메트릭을 모니터링할 수 있는 웹 플랫폼을 설계하고 개발했습니다.
  \item 공식 C++ API와 ROS 패키지를 포함한 회사 오픈 소스 리포지토리의 개발, 릴리스 주기 및 지속적인 유지보수를 주도하고 있습니다.
\end{cvbullets}

\cvrole{마키나락스}{2020/09--2025/08}{로보틱스 \& 머신러닝 리서치 엔지니어}{}
\begin{cvbullets}
  \item 수백 대의 로봇이 배치된 점용접 조립 라인을 위한 로봇 프로그램 자동 생성 시스템의 소프트웨어 아키텍처를 주도했습니다.
  \item 산업용 로봇 프로그래밍 소요 시간을 약 6주에서 3일로 단축했습니다.
  \item ROS, MoveIt, OMPL 및 MongoDB를 활용해 고도로 병렬화된 C++ 플래닝 알고리즘을 개발하여 작업점 검증, 궤적 생성 및 평가, 작업 분배, 충돌을 고려한 로봇 간 협조를 구현했습니다.
  \item 연산 집약적인 플래닝 워크로드를 위한 Kubernetes/Docker 기반 확장형 배포 인프라를 구축·운영하고, 빌드·테스트·패키징·배포를 위한 Linux/GitHub Actions CI/CD를 유지보수했습니다.
  \item PyTorch, Ray RLlib, ML-Agents 및 Unity3D를 사용해 모방학습과 강화학습 실험 및 시뮬레이션 개발을 수행했습니다.
  \item 모바일 엣지 디바이스 배포를 위해 LLM 아키텍처를 모델 수준에서 최적화했습니다.
  \item PyTorch/ONNX 기반 그래프 최적화, 양자화 및 반복 가능한 성능 평가 워크플로를 구축했으며 업무 범위는 모델 및 그래프 수준에 한정되었습니다.
\end{cvbullets}

\cvrole{Moringa Digital / Portal de Compras Públicas}{2016/08--2017/07}{소프트웨어 개발자}{}
\begin{cvbullets}
  \item 공공 조달 플랫폼과 고객 ERP 시스템 간 연동을 자동화하는 Node.js 서비스를 개발했습니다. 이 연동은 플랫폼 운영 시간을 줄이고 회사의 주요 영업 포인트가 되었으며, 웹사이트 백엔드 개발과 고객 기술팀 협업도 수행했습니다.
\end{cvbullets}

\cvmainsection{학력 \& 연구}
\cvrecord{전자전기컴퓨터공학 석사}{2018--2020}{성균관대학교, GPA 4.4/4.5. 시맨틱 지식 프레임워크 연구에 참여하고 ROS와 Gazebo 기반의 자동 심적 시뮬레이션 시스템을 개발했으며, 강화학습 기반 내비게이션 및 하이브리드 경로 계획 기법을 연구했습니다.}
\cvrecord{제어 및 자동화공학 학사}{2011--2016}{University of Brasilia, GPA 4.1/5.0. 듀얼 쿼터니언 기반 협동 조작 방법을 도출하고 NAO와 V-REP에서 제어 전략을 구현했으며 OpenCV 기반 로봇 간 시각 통신을 개발했습니다.}
\cvrecord{컴퓨터공학 교환 과정}{2014--2015}{성균관대학교, GPA 4.1/4.5}

\cvmainsection{기타 활동 \& 봉사}
\cvrecord{UnBeatables 휴머노이드 로봇 축구팀}{2015--2016}{팀장으로서 RoboCup과 LARC 참가를 이끌고 멀티스레드 C++/Python 아키텍처, UDP 통신, 행동 관리 및 휴머노이드 모션 통합을 개발했습니다. 학교, 어린이 병원 및 과학 행사에서 로봇 시연과 STEM 진로 홍보 활동을 진행했습니다.}
\cvrecord{Electron Project}{2013--2014}{브라질리아 공립학교 학생들에게 전자공학과 프로그래밍을 가르친 자원봉사 교사}

\cvmainsection{오픈 소스}
\cvrecord{\href{https://github.com/yurirocha15/mcp-cpp-sdk}{mcp-cpp-sdk}}{}{Model Context Protocol용 C++20 SDK}

\cvmainsection{자격, 수상 \& 언어}
\cvrecord{\href{https://www.credly.com/badges/bb876a9e-74e1-475d-b753-c43f2675c9ee}{Certified Kubernetes Administrator}}{}{LF-eivj0wxw1q}
\cvrecord{\href{https://www.credly.com/badges/f1d71150-7469-4e33-9745-758d7256fa35}{Certified Kubernetes Security Specialist}}{}{LF-nyf5h9e9v9}
\cvrecord{\href{https://forums.developer.nvidia.com/t/the-results-are-in-meet-the-nvidia-cosmos-cookoff-winners-see-them-live-on-april-16/366130}{NVIDIA Cosmos Cookoff 우승}}{2026}{\href{https://github.com/doosan-robotics/explainable-palletizer}{Explainable mixed palletizing --- 공식 코드 공개}}
\cvrecord{정부초청장학생 학업성취상}{2019}{정부초청외국인장학사업(GKS)}
\cvrecord{정부초청외국인대학원장학생}{2017--2020}{성균관대학교 대학원 장학금}
\cvrecord{언어}{}{포르투갈어: 모국어 \textbullet\ 영어: 유창, TOEFL 115/120 \textbullet\ 한국어: 유창, TOPIK 5급 \textbullet\ 프랑스어: 중급}

\cvmainsection{연구 실적}
\cvrecord{\href{https://www.mdpi.com/2076-3417/10/9/3219/pdf}{Autonomous Navigation Framework for Intelligent Robots Based on a Semantic Environment Modeling}}{2020}{S.-H. Joo*, S. Manzoor*, Y. G. Rocha, S.-H. Bae, K.-H. Lee, T.-Y. Kuc, M. Kim. Applied Sciences, 10:3219.}
\cvrecord{\href{https://www.yurirocha.com/assets/Yuri_Master_Thesis.pdf}{Mental Simulation for Autonomous Learning and Planning Using an Ontology-Based Modeling Framework}}{2020}{Y. G. Rocha. 성균관대학교 석사학위 논문.}
\cvrecord{\href{https://www.yurirocha.com/assets/Mental_Simulation_IROS_2019.pdf}{Mental Simulation for Autonomous Learning and Planning Based on Triplet Ontological Semantic Model}}{2019}{Y. G. Rocha, T.-Y. Kuc. CEUR Workshop Proceedings, 2487:65--73.}
\cvrecord{\href{https://www.yurirocha.com/assets/Automatic_Generation_ICCAS2019.pdf}{Automatic Generation of a Simulated Robot from an Ontology-Based Semantic Description}}{2019}{Y. G. Rocha, S.-H. Joo, E.-J. Kim, T.-Y. Kuc. ICCAS, 1340--1343.}
\cvrecord{\href{https://www.yurirocha.com/assets/Design-of-singularity-robust-and-task-priority-primitive-controllers_CCTA_2017.pdf}{Design of Singularity-Robust and Task-Priority Primitive Controllers for Cooperative Manipulation}}{2017}{C. M. de Farias*, Y. G. Rocha*, L. F. C. Figueredo, M. C. Bernardes. IEEE CCTA, 740--745.}

\cvmainsection{등록 특허}
\cvrecord{로봇 공정을 위한 프로그램을 생성하는 방법}{\href{https://patents.google.com/patent/KR102590491B1}{KR102590491B1}}{}
\cvrecord{로봇의 유효 작업 지점을 결정하기 위한 방법}{\href{https://patents.google.com/patent/KR102626109B1}{KR102626109B1}}{}
\cvrecord{산업용 로봇의 작업 경로 생성 방법}{\href{https://patents.google.com/patent/KR102629021B1}{KR102629021B1}}{}
\cvrecord{작업 수행 로봇의 작업 경로의 길이를 계산하는 방법}{\href{https://patents.google.com/patent/KR102660168B1}{KR102660168B1}, \href{https://patents.google.com/patent/KR102638245B1}{KR102638245B1}}{}
\cvrecord{복수의 작업 수행 로봇들에 작업 지점들을 분배하기 위한 방법}{\href{https://patents.google.com/patent/KR102614099B1}{KR102614099B1}, \href{https://patents.google.com/patent/US12093832B2}{US12093832B2}}{}
\cvrecord{복수의 작업 수행 로봇들 간 충돌을 방지하기 위한 방법}{\href{https://patents.google.com/patent/US12049013B1}{US12049013B1}}{}
\end{document}
